 \section{Presentation Plan}

\subsection*{Intro / project choice (30sec)}
Hi, my name is Henry and welcome to my EPQ Project

\textbf{Dramatic Drone Shot logo reveal with music, pans up from me to reveal logo}

Neuronote.

For my EPQ I wanted to explore the intersection of human learning and technology \textbf{animation of the two icons combining} and chose to implement my research by creating a digital note-taking platform from the ground up \textbf{screenshot of website moves up from ground with rubble effect} which helps users to learn more effectively using a variety of different input methods and techniques but how did I get here?,

\textbf{Whoosh back in time sound effect}

\textbf{Flies through O into gradient background slides}

\subsection*{Project Choice (1min)}
To start my EPQ Journey I spent some time completing a personal audit to explore my preexisting interests, skills and passions \textbf{Audit on screen}

I found this process extremely valuable, as it distilled what I've enjoyed most in life into a concise document. From this, I realized I wanted to pursue a project that combined my academic research skills with a personal passion.

To organize my thoughts on various topics, I created a few mind maps that grouped my ideas into different categories, \textbf{mindmaps}

With a growing interest in computer science, I initially leaned toward a dissertation-style project (similar to my previous HPQ) that would allow me to explore computational data and its impact on human thought. However, I didn't want to commit immediately and planned to spend a couple of weeks selecting my final project. \textbf{clock spin animation}

At this point I was fortunate to spend some time on a work placement with a phenomenal Cambridge Technology startup who are developing revolutionary software to help companies organise their data and workflows in a very simple, intuitive yet powerful way. \textbf{Kinabase logo and website viusal} Having spent some time with the team and seeing the real impact of their software on day-to-day operations,   I decided to change my focus away from a dissertation towards a more practical, Artifact project. This would allow me to create a tool that could genuinely help people in their daily lives while also strengthening my future university applications. 

While searching for theme inspiration I reflected on my Mind-Map \textbf{reshow mindmap} and was drawn to the psychology of learning \textbf{highlight this area}. As a student I have often felt as if the platforms I had been using didn't 'click' with me. \textbf{show logos of a few different noteettaking apps, onenote, notability and obsidian} The opportunity to build my own tool, grounded in scientific research and tailored to what my peers wanted, was therefore very appealing.”

\subsection*{Planning (1min)}
\textbf{Transition stinger}

Having chosen an ambitious project subject I knew that careful, coordinated planning would be crucial to the success of my artifact.
 
I began by laying out my goals within a project choice document. I was already aware that scope creep could quickly hinder my progress, and so I made an Eisenhower matrix of targets \textbf{Show matrix}, ranging from crucial high impact goals to lesser used but nonetheless nice to have features. With my project now scoped I developed my first timeline. \textbf{Timeline}
 
Having not previously worked on a project this large before, this first plan was very high level giving me flexibility to tighten down with much more information once I had researched and learnt about the technicalities in detail. I originally started with a Gantt chart having researched its powerful benefits \textbf{show Gant chart} however I soon reflected on the timeline I'd produced and found that my project was better split into three discrete themes: Human learning, user experience and technical artifact development. \textbf{Make three pillars on screen}

I found that this separation of concerns suited a more traditional 'side by side' timeline \textbf{animate timeline below}  as I could quickly cross-reference multiple goals set for the same week in a nice visible way whilst also focusing on different themes when required. With a huge amount to learn ahead of me I blocked out weekly slots in my diary that I could use to focus on my project. This would later be one of the most valuable decisions I made early on in  as it allowed me to spend time focused on my goals without the distractions of other subjects.

As my project progressed I continued to update my plans. On a few occasions I found myself ahead of schedule allowing me to spend time on some of my lesser priority features (in accordance with my Eisenhower matrix). \textbf{Montage of coding the chemistry formula feature}

Some time periods (such as EPQ Day or Holidays) enabled me to spend much larger blocks of time on my project, and so I also made more detailed task sheets for these to maximize my productivity. \textbf{show EPQ day plan}


\subsection*{Research (1min)}
\textbf{Transition stinger}
In line with my theme split whilst planning I also tackled my research as a triad approach of Technical (for software development), Theoretical (human learning) and Tactical (user experience) research.

I found that these different areas each found themselves more naturally aligned with different research methods. For example: technical software favoured more modern, fast moving media such as web-based documentation, YouTube tutorials and podcasts, whilst research on human learning saw me study a wide range of scientific articles dating back to the 1900s and beyond. \textbf{show different icon representing each type of research}

Interestingly user experience research leaned most strongly on my interactions with peers and primary users as it was a largely subjective process. This primary research comprised of an initial survey on peoples perception of existing note-taking methodologies, with later interviews, surveys and self-guided research on existing platforms helping me further improve and reflect on the positives and improvements points within my project. \textbf{cover screen with variety of survey and interview screenshots}


Overall I used over 50 sources \textbf{show bibliographys} across a wide variety of categories. My favourite source is a report written by Daniel Oppenheimer titled "The Pen Is Mightier Than the Keyboard" \textbf{show report page on journal website} which explored the advantages handwriting notes had on information recollection among students. This paper guided me to build a written input type within my platform which was hugely popular with my primary users. This paper was particularly useful as it was heavily peer-reviewed with supplementary research conducted by Kayla Morehead and others titled: "How Much Mightier Is the Pen than the     for Note-Taking". \textbf{show other linked paper }

From the outset of my research I was careful to keep not just my detailed notes on the source but also a reference within my source analysis table so that I could easily refer to it later on. \textbf{show source analysis table} 

Whilst this worked well early on, as my research grew I found it increasingly difficult to keep track of which source correlated to which research, and so I switched to using Zotero \textbf{show zotero logo}, a free and open source reference management software. Zotero allowed me to easily collate my sources and even generate citations in a variety of formats \textbf{screen recording of using zotero}. This was a huge time saver and I would recommend it to anyone who is looking to conduct extended research in the future. I also found that using Zotero helped me to properly credit libraries and research I used and avoid plagiarism, as I could easily see which sources I had used in my research and how they related to my project.

\subsection*{Skill Development (1min)}
Throughout my project I have been a  ble to develop a large range of skills throughout a wide range of categories. To help me later reflect on my skill progression I completed an initial skills audit at the start of my EPQ which ranked 21 key skills which I identified as crucial to my progression and sucess. \textbf{show skills audit}

As my skills developed I completed Additional audits, and it was hugely rewarding to visually see my personal development throughout the project \textbf{show three skills audits on screen}. In addition to these audits I recorded a robust entry in my diary log for every week in the EP program which allowed me to reflect on my progression within each individual skill category I assessed \textbf{fly through of many diary entries slanted on left with a static one on right that is legible}. I then produced an overall reflection for each which highlights my progression throughout the project in that skillset. \textbf{same again but fly on right and static left}

During the project I found that I learn bes t by interacting with the material in question in a hands-on way, and so I made heavy use of prototyping and iterative design as I was learning new skills such as User interface design, Website development and server hosting . \textbf{show a variety of prototypes on screen}

The mistakes I made here helped guide me to rethink my processes in new, refined ways which helped elevate my artifact and hence the skills I was developing. Throughout my project I recorded each of these prototypes and iterations of which a few you can see on screen now.

Some of these prototypes were very simple, testing basic features such as users logging in, some were slightly stranger ideas that didn't meet the cut and some continued throughout my project and found themselves embedded in my final artifact.

Throughout this time peer feedback was crucial to help me understand any friction points and steer me in better directions

as I acquired more skills, these prototypes grew and grew in complexity and by the beginning of March I was ready to build my final artifact. \textbf{text "Final Artifact" appears on screen}


\subsection*{Artifact Development (1min)}
Combining all of the technical skills I had developed, the understanding of human learning I had and lots of code I'd written in prototypes I got to work building. \textbf{coding montage}


Whilst my research was very thorough I knew that I would likely hit some bumps on the way and made sure my project plan was flexible enough to allow me to adapt to these challenges. One notable example of this was during a conversation with a less technically inclined user who was struggling to understand the link between notes on the graphs and the tags they were assigned. Whilst I was able to simply explain my concept to them and how it differed to more traditional platforms I realised that this would likely be a common friction point, and so I created an interactive tutorial  to help users understand the concept of tags and how they can be used to link notes together. \textbf{tutorial video}

I aimed my project to feel 'commerical grade' and so I even dived into some Legal skills as I wrote my privacy and terms of use policies guided by my research \textbf{show policy section on webssite (login and settings)}. These elements are often forgotten in student projects, however I wanted to challenge myself and further grow my skill set. Equally, I also made sure to include a section on accessibility within my project, as I wanted to ensure that my platform was usable by as many people as possible. This included making sure that my platform was compatible with screen readers and other assistive technologies in line with government guidance , as well as ensuring that it was easy to navigate for users with differing and diverse needs. \textbf{accessibility section from notes}

\subsection*{Evaluation (1min)}
Having built a platform focused on improving student learning the comments from my peers would serve as crucial feedback, and they didn't disappoint. On the 27th of June I launched NeuroNote online. \textbf{Neuronote home screen} I emailed my waiting list, messaged my friends and family and posted on social media to announce the launch of my platform. \textbf{show messages I sent and reddit post} After a nervous few minutes as people signed up and inspected my work. I was overwhelmed by the positive response and reception I received. Over 100 people signed up to use my platform within the first week and many are still regularly logging in \textbf{show Cloudflare viewing stats}. From the feedback forms completed it is clear that users believe my platform gives them a space to better organise their notes, thoughts and ideas, exactly whay i had set out to achieve

When I look back and use my project in my day-to-day life I am incredibly proud of what I have achieved. I have been able to create a platform that is not only functional but also visually appealing and easy to use. I have also achieved all of this whilst balancing my other A-levels and commitments, 

That isn't to say however that everything went perfectly, along the way I encountered many challenges and setbacks, but I was able to overcome them by using the skills I had developed throughout the project. For example my Server running from my house suffered a power cut on launch day \textbf{show uptime robot}, fortunately, due to monitoring I had set up and remote access I had configured  \textbf{show terminal app on phone and tailscale} I knew in a matter of seconds and was able to redirect traffic to the cloud whilst I restored my house connection resulting in only 2 minutes of downtime since going live. If I were to do it all again my main change would be to be less set on achieving perfection within every area of the project, especially whilst prototyping, as I found myself spending too long on small details which didn't really matter in the grand scheme of things, when the end outcome wouldn't even make it into my project. 

   Secondly I would be more formal in my primary research creation, Whilst I did make use of forms, feedback portals and interviews, I had so many incredible conversations face to face and found myself wishing I had recorded them for later reference. Lastly I would have liked to have dedicated more time on the final presentation of my project, I feel that I could have created a more polished portfolio by creating some adverts, instructional videos and other marketing materials to help promote my platform on social media.

Despite this the future of my project is very exciting, within the coming months I aim to host my code publically on Github so that anyone can use (and maybe even contribute to) my project. Having built the platform I always wished I had, NeuroNote will continue to be my go-to platform for all my notes throughout Upper Sixth, university and beyond which I strongly believe will improve my academic performance. Lastly I did make a short product announcement trailer demonstrating many of the benefits of NeuroNote which I hope you will enjoy:
